problem:  
Let \((M^n, g)\) be a complete, noncompact Riemannian manifold with nonnegative sectional curvature and **Euclidean volume growth**, i.e.  
\[
\lim_{r\to \infty} \frac{\operatorname{Vol}(B(p, r))}{r^n} = v_0 > 0,
\]
for some (and hence any) base point \(p\in M\).  
Assume the curvature tensor satisfies
\[
|Rm|(x) \le C\, r(x)^{-(2+\alpha)}
\]
for some constants \(C>0\) and \(\alpha>0\), where \(r(x)=d(p, x)\).  

**Conjecture:**  
There exists a dimensional constant \(\alpha_0 > 0\) such that if \(\alpha > \alpha_0\), then \(M\) has a **unique** tangent cone at infinity, and this tangent cone is isometric to \(\mathbb{R}^n\).  

Equivalently, sufficiently fast curvature decay (beyond a critical rate \(\alpha_0\)) enforces *asymptotic flatness with unique Euclidean tangent cone*.  

Moreover, determine whether the value of \(\alpha_0\) can be characterized geometrically—e.g., via the rate at which the volume ratio  
\[
\frac{\operatorname{Vol}(B(p, r))}{\omega_n r^n}
\]
approaches its limit \(v_0\), or through the vanishing rate of the Abresch–Gromoll excess function.  

Why is it a "good" problem:  
This conjecture precisely formulates a quantitative rigidity problem connecting **curvature decay**, **volume growth**, and **asymptotic structure** for manifolds with nonnegative sectional curvature. It offers a natural sharpening of soul and splitting theories by asking how *fast curvature must decay* to force the global geometry at infinity to stabilize to Euclidean space.  

The problem is mathematically well-posed and open: while the existence of some tangent cone at infinity follows by Gromov compactness, no known theorem identifies quantitative curvature-decay thresholds that guarantee **uniqueness** or **Euclideanity** of that cone. By incorporating the precise decay exponent \(2+\alpha\) (the scaling of the curvature tensor), the problem aligns with geometric analysis and scaling invariance principles.  

It bridges multiple domains—comparison geometry (nonnegative curvature and excess estimates), metric geometry (tangent cones and Gromov–Hausdorff limits), and geometric analysis (quantitative decay estimates). The statement is both simple and deep: it asks, in essence, whether “fast-decaying curvature forces Euclidean infinity.” That simplicity of formulation yet profound depth of consequence makes it a *good mathematical problem*.